\section{Executive Investment Recommendation}

Synthesizing empirical findings from data acquisition, exploratory return dynamics, mean-variance optimization, and downside risk analytics, this section formulates quantitative investment recommendations for Group 3 (Charlie).

\subsection{Core-Satellite Portfolio Architecture}

We recommend establishing a \textbf{Core-Satellite Allocation Structure (60 percent Max Sharpe / 40 percent GMV tilt)}.

\begin{itemize}
  \item \textbf{Core Stabilizer (40 percent GMV Allocation):} Anchored in domestic defensive equities (\texttt{TEL}: 15.93 percent, \texttt{AEV}: 12.74 percent, \texttt{MER}: 9.92 percent), providing low return correlation relative to global tech stocks, dividend income, and downside capital protection.
  \item \textbf{Satellite Growth Engine (60 percent Max Sharpe Allocation):} Captures high-momentum equity capital appreciation via \texttt{NVDA} (35.02 percent), \texttt{META} (13.56 percent), and domestic power utility \texttt{MER} (11.42 percent).
\end{itemize}

This hybrid structure achieves a balanced quantitative profile: an annualized target expected return of approximately 33.5 percent, an annualized volatility of 21.5 percent, and a Sharpe ratio of approximately 1.05 relative to the 5.25 percent Philippine 91-Day T-Bill benchmark.

\subsection{Execution Governance and Rebalancing Protocol}

To maintain target risk exposures and capture systematic rebalancing gains, institutional execution should enforce three operational risk controls:

\begin{enumerate}
  \item \textbf{Mandatory Monthly Rebalancing Calendar:} Empirical backtesting reveals that un-rebalanced Buy-and-Hold strategies experience extreme asset weight drift, allowing volatile assets like NVIDIA and Bitcoin to dominate portfolio risk. Monthly rebalancing systematically trims outperforming assets at elevated valuations and reallocates capital into core defensive positions, capturing a rebalancing premium ($+0.5 \sigma_p^2 (1 - \bar{\rho})$).
  \item \textbf{Daily Value at Risk Budgeting Limit:} Enforce a maximum daily 95 percent Historical VaR risk budget ceiling of \textbf{2.20 percent}. If daily portfolio drawdown breaches \textbf{-3.50 percent} (95 percent Expected Shortfall threshold), trigger a dynamic de-risking protocol that temporarily reallocates satellite growth exposure into Philippine 91-Day Treasury Bills ($R_f = 5.25\%$).
  \item \textbf{Liquidity Reserve Buffer:} Maintain a 5.0 percent cash or 91-Day T-Bill liquidity reserve buffer to absorb market volatility and facilitate rebalancing execution without forced asset liquidations during market stress.
\end{enumerate}
